% ===================================================================
%  ТАБЛИЦЯ № 1 — Школа. — Ліцей. — Клас.
%
%  Ceci n'est PAS une transcription : c'est une traduction, faite en
%  2026, en ukrainien d'aujourd'hui. D'ou trois differences avec
%  texto/io et texto/fr, et trois seulement :
%
%   * pas de \nl ni de \cc — la traduction ne suit aucune ligne
%     imprimee, et les lignes de ce fichier ne sont que du confort ;
%   * pas de \begin{VUpage} — elle n'a ni feuillet ni folio ;
%   * le gras se pose sur le TERME seul, ponctuation dehors.
%
%  Les cles « %%K » sont celles de texto/io, macro pour macro pour
%  l'apparat, et les renvois \textsuperscript{(n)} visent les MEMES
%  objets de la planche.
%
%  HUITIEME COLONNE DE LA SERIE DES DIX-SEPT, et la premiere de ce
%  groupe qui s'ecrive en cyrillique. La colonne russe est deja
%  servie ; elle sert donc de temoin, comme l'espagnol a servi au
%  catalan et le catalan a l'occitan. TEMOIN, ET NON MODELE : c'est
%  precisement la lecon des deux colonnes precedentes. Le castillanisme
%  passe en catalan par l'oreille, le catalanisme en occitan de meme ;
%  le russisme passerait ici par la meme porte, et pour la meme
%  raison — la ressemblance dispense de verifier.
%
%  ET LES DEUX LEXIQUES SCOLAIRES DIFFERENT PLUS QU'ON NE CROIT. Le
%  russe « шкаф » est masculin, l'ukrainien « шафа » est feminin ; le
%  russe dit « мел », l'ukrainien « крейда », qui est feminin aussi ;
%  « тетрадь » est feminin en russe, « зошит » masculin en ukrainien ;
%  la recreation est « перемена » la-bas et « перерва » ici ;
%  l'horloge est « часы », un pluriel, la-bas, et « годинник », un
%  singulier, ici. Cinq mots pris au hasard du premier tableau, cinq
%  differences : la traduction se fait donc mot a mot sur l'ido, et le
%  russe ne sert qu'a verifier apres coup quel objet de la planche est
%  vise.
%
%  UNE LANGUE A CAS, comme le russe. Le renvoi suit son substantif, et
%  le substantif se decline : le gras porte donc la forme FLECHIE,
%  celle que la phrase demande — « на \VUgras{лавах} », non le
%  nominatif.
%
%  SON GENITIF SE POSTPOSE, comme celui de l'ido : « двері старого
%  замку », non « старого замку двері ». On attend donc zero inversion,
%  comme en catalan et en occitan, et pour la meme raison de fond :
%  l'ordre du complement du nom.
%
%  LE NOM DE L'EDITEUR SE PRONONCE « del-MASS », et l'ukrainien, qui
%  decline les noms d'homme termines par une consonne, demande le
%  genitif : « таблиць Дельмаса ».
% ===================================================================

%%K t01-apar-0 apar
\VUsaut{2.0mm}
\VUcentre{12.6pt}{120}{ПОЯСНЮВАЛЬНА КНИЖЕЧКА}
\VUsaut{3.2mm}
\VUcentre{8.4pt}{160}{ДО}
\VUsaut{2.6mm}
\VUtitre{15.6pt}{0}{допоміжних таблиць Дельмаса}
\VUsaut{3.4mm}
\VUornamento{9pt}
\VUsaut{5.0mm}
\VUcentre{13.2pt}{90}{ПЕРША СЕРІЯ}
\VUsaut{3.0mm}
\VUfilet{16mm}
\VUsaut{5.6mm}

%%K t01-apar-1 apar
\VUtitre{15.0pt}{60}{ТАБЛИЦЯ № 1}
\VUsaut{1.4mm}
\VUtitre{11.4pt}{0}{Школа. --- Ліцей. --- Клас.}
\VUsaut{2.2mm}
\VUfilet{20mm}
\VUsaut{6.4mm}

%%K t01-tit-1 sub
\VUpk{10.2pt}{40}{Загальний опис.}
\VUsaut{3.0mm}

%%K t01-01-1 p
1. --- На цій таблиці зображено \VUgras{клас} у \VUgras{ліцеї}. У цьому класі
вісім \VUgras{столів} \textsuperscript{(18)} і вісім \VUgras{лав}
\textsuperscript{(32)}. На кожній лаві сидять три учні. \VUgras{Учитель}
\textsuperscript{(7)} сидить на \VUgras{стільці} \textsuperscript{(8)} за своєю
\VUgras{кафедрою} \textsuperscript{(6)}.

%%K t01-02-1 p
2. --- Ліворуч від кафедри стоїть заскле́на \VUgras{шафа} \textsuperscript{(15)}.
Праворуч --- \VUgras{класна дошка} \textsuperscript{(1)} з \VUgras{ганчіркою}
\textsuperscript{(2)} і \VUgras{крейдою} \textsuperscript{(3)}.

%%K t01-02-2 p
Над кафедрою висять \VUgras{настінні таблиці} \textsuperscript{(9, 11, 12)} і
\VUgras{мапа} \textsuperscript{(10)}.

%%K t01-03-1 p
3. --- Не всі учні на своїх \VUgras{місцях} \textsuperscript{(17)}. Іван
\textsuperscript{(4)} стоїть коло дошки; він тримає ганчірку і стирає
\VUgras{геометричні фігури}.

%%K t01-03-2 p
Петро стоїть коло кафедри; він збирає \VUgras{письмові роботи}
\textsuperscript{(5)} і несе їх учителеві.

%%K t01-04-1 p
4. --- \VUgras{Цей} учитель --- викладач \VUgras{нових мов}. Він учить учнів
говорити, читати й писати чужими мовами \VUgras{(німецькою, англійською,
іспанською, італійською, російською} або \VUgras{французькою)}.

%%K t01-05-1 p
5. --- Клас дуже великий і дуже світлий. У глибині широке \VUgras{вікно} з
двома \VUgras{кватирками} \textsuperscript{(78)} і \VUgras{заскле́ні двері}, щоб
його провітрювати й освітлювати.

%%K t01-05-2 p
У цьому класі не холодно. Посередині, щоб його обігрівати, стоїть дуже добра
\VUgras{піч} \textsuperscript{(46)} зі своїм \VUgras{димарем}
\textsuperscript{(45)}.

%%K t01-05-3 p
До перегородки прибиті \VUgras{вішалки} \textsuperscript{(58)}; на них висять
капелюхи й пальта учнів.

%%K t01-tit-2 sub
\VUsaut{5.2mm}
\VUpk{10.2pt}{40}{Подвір'я.}
\VUsaut{3.6mm}

%%K t01-06-1 p
6. --- \VUgras{Годинник} \textsuperscript{(63)} показує п'ять хвилин на десяту.

%%K t01-06-2 p
\VUgras{Перерва} \textsuperscript{(76)} щойно скінчилася, і урок починається
знову. Перерва минає на \VUgras{подвір'ї} \textsuperscript{(75)}. Ми бачимо
кількох учнів, що граються. \VUgras{Наглядач} \textsuperscript{(74)} стежить за
ними.

%%K t01-07-1 p
7. --- На подвір'ї ми бачимо й інших людей: \VUgras{інспектора}
\textsuperscript{(73)}, \VUgras{директора} \textsuperscript{(71)} і
\VUgras{завуча} \textsuperscript{(72)}.

%%K t01-07-2 p
Інспектор перевіряє школи, директор керує закладом, завуч стежить за працею
учнів.

%%K t01-07-3 p
Четвертий --- \VUgras{сторож} \textsuperscript{(77)}. Він несе під пахвою
журнал, щоб відзначити відсутніх.

%%K t01-08-1 p
8. --- У глибині подвір'я стоять \VUgras{гімнастичні прилади}
\textsuperscript{(70)} і майстерня \VUgras{ручної праці} \textsuperscript{(69)}.

%%K t01-08-2 p
Праворуч на таблиці ледве видно \VUgras{класи} \VUgras{музики} і
\VUgras{малювання}.

%%K t01-noto-1 noto
\VUnotes{143.2mm}{%
(*) Для \textit{імен} теж радимо передавати їх у латинській формі. Тільки так
можна уникнути незліченних різночитань. Та й як вибрати між \textit{Giovanni,
Jean, Johann, Jan, Hans, John, Іван} та іншими? Не створювати ж окремого
словника імен? Візьмімо латинський називний відмінок і скажімо:
\VUgras{Ioannes, Ioanna; Iulius, Iulia; Iakobus; Andreas; Lukas.} (Повна
граматика).%
}

%%K t01-tit-3 sub
\VUpk{10.2pt}{40}{Докладний опис.}
\VUsaut{2.4mm}

%%K t01-tit-4 sub
\VUpk{10.2pt}{40}{Добрі учні.}
\VUsaut{3.0mm}

%%K t01-09-1 p
9. --- Учні на першій лаві праворуч від кафедри --- \VUgras{Генріх} (*),
\VUgras{Стефан} і \VUgras{Карл}.

%%K t01-09-2 p
Генріх дуже уважний і старанний; він слухає вчителя. На столі в нього
\VUgras{зошит} \textsuperscript{(28)}, \VUgras{книжка} \textsuperscript{(29)},
\VUgras{словник} \textsuperscript{(30)} і \VUgras{пенал} \textsuperscript{(31)}.
У його книжці багато \VUgras{сторінок}; у кожному розділі кілька
\VUgras{абзаців}, а в кожному \VUgras{рядку} багато \VUgras{слів}.

%%K t01-09-4 p
Його сусід Стефан \textsuperscript{(24)} ретельний. Він записує в зошит урок,
заданий на наступний раз. У нього прегарний \VUgras{почерк}. Книжка його
закрита. Видно тільки \VUgras{оправу} \textsuperscript{(27)}. Поруч із ним
лежить його \VUgras{циркуль} \textsuperscript{(26)}.

%%K t01-09-5 p
Карл \textsuperscript{(23)} тримає книжку в руці; він розгортає її, щоб іще раз
повторити урок. Як і в кожного учня, перед ним стоїть \VUgras{чорнильниця}
\textsuperscript{(25)}. Він слухняний.

%%K t01-10-1 p
10. --- За наступним столом сидять \VUgras{Людовик, Антон} і \VUgras{Павло}.
Людовик відповідає свій \VUgras{урок} \textsuperscript{(22)} і відповідає на
\VUgras{запитання} вчителя. Антон \textsuperscript{(21)} дуже неуважний; інколи
вчитель навіть карає його. Павло \textsuperscript{(20)} --- добрий товариш,
привітний і послужливий. Він дуже уважний.

%%K t01-tit-5 sub
\VUsaut{4.2mm}
\VUpk{10.2pt}{40}{Погані учні.}
\VUsaut{3.2mm}

%%K t01-11-1 p
11. --- Позаду Генріха сидить \VUgras{Георгій} \textsuperscript{(38)}. Він не
злий хлопець, але \VUgras{балакучий}, неуважний і непосидючий. Його
\VUgras{легковажність} і \VUgras{неслухняність} вносять \VUgras{безлад} в урок,
і він часто заслуговує на суворе \VUgras{зауваження}. Його \VUgras{чернетка}
\textsuperscript{(35)} брудна, він ставить на ній \VUgras{ляпки} своїм
\VUgras{пером} \textsuperscript{(34)}, через що постійно береться за
\VUgras{гумку} \textsuperscript{(36)}; його \VUgras{ранець}
\textsuperscript{(33)} подертий, такий він неохайний. Навіщо він приносить свою
\VUgras{ручку} \textsuperscript{(37)} до класу нових мов?

%%K t01-11-2 p
Його сусід Едмунд \textsuperscript{(39)} його не любить. Щоправда, він трохи
лінивий, але тихий і спокійний. Він не балакає; тільки замість того, щоб
слухати, воліє читати \VUgras{оповідання} у своїй \VUgras{хрестоматії}.

%%K t01-11-3 p
Третій учень на цій лаві --- \VUgras{Людвіг} \textsuperscript{(40)}. Він
старанний. Він пише на білому \VUgras{аркуші паперу}. Перед собою він поклав
\VUgras{промокальний папір} \textsuperscript{(41)}.

%%K t01-12-1 p
12. --- Учні на другій лаві, ліворуч від учителя, зовсім малі. Перший, малий
\VUgras{Еміль}, ще не вміє добре писати; він приносить свою \VUgras{грифельну
дошку} \textsuperscript{(42)}, свій \VUgras{грифель} і свою \VUgras{губку}
\textsuperscript{(43)}.

%%K t01-12-2 p
Поруч із ним \VUgras{Август} і \VUgras{Бернард} \textsuperscript{(44)}. Останній
добре працює, але \VUgras{вигляд} у нього дуже неохайний.

%%K t01-13-1 p
13. --- Позаду них сидять троє \VUgras{ледарів}. \VUgras{Альберт} не вчить
уроків. \VUgras{Яків} \textsuperscript{(47)} спить замість того, щоб слухати.
Його \VUgras{лінощі} накликають на нього \VUgras{догани} і навіть
\VUgras{покарання}. Нарешті, \VUgras{Християн} \textsuperscript{(48)} надто
легковажний. Він погано \VUgras{поводиться} і не намагається виправитися.

%%K t01-14-1 p
14. --- Їхні сусіди, навпаки, поводяться добре. \VUgras{Ернест}
\textsuperscript{(51)} любить лад і точність; він завжди користується своєю
\VUgras{лінійкою} \textsuperscript{(50)} і своїм \VUgras{олівцем}
\textsuperscript{(49)}. Він \VUgras{екстерн}.

%%K t01-14-2 p
\VUgras{Франциск} \textsuperscript{(52)} --- \VUgras{пансіонер}; він носить
однострій, їсть і спить у закладі. Він добрий \VUgras{товариш}.

%%K t01-14-3 p
Моріц загострює свій \VUgras{олівець} \VUgras{складаним ножиком}
\textsuperscript{(56)}; він дуже ощадливий.

%%K t01-14-4 p
Він приносить усі свої книжки, свою \VUgras{граматику} \textsuperscript{(55)},
свій \VUgras{записник} \textsuperscript{(54)} і навіть \VUgras{блокнот}
\textsuperscript{(53)}. Його \VUgras{портфель} \textsuperscript{(57)} лежить на
підлозі позаду нього.

%%K t01-14-5 p
Два столи в глибині класу зайняті \VUgras{добрими учнями}
\textsuperscript{(19)}.

%%K t01-tit-6 sub
\VUsaut{4.6mm}
\VUpk{10.2pt}{40}{Малювання і музика.}
\VUsaut{3.4mm}

%%K t01-15-1 p
15. --- У \VUgras{класі малювання} на стіні висять гарні \VUgras{зразки}, а зі
стелі звисає \VUgras{лампа} \textsuperscript{(62)} зі своїм \VUgras{абажуром}.
\VUgras{Учитель} \textsuperscript{(60)} дуже терплячий; він виправляє
\VUgras{рисунки} \textsuperscript{(61)} учнів і вчить їх рисувати
\VUgras{погруддя} \textsuperscript{(59)} з натури.

%%K t01-16-1 p
16. --- \VUgras{Клас музики} здається дуже цікавим. Там учаться читати
\VUgras{ноти}, співати \VUgras{щаблі гами} \textsuperscript{(67)}, написані на
\VUgras{нотному стані} (до, ре, мі, фа, соль, ля, сі\,=\,C. D. E. F. G. A. B.),
знати \VUgras{ключі} і співати чарівні \VUgras{мелодії}. Ноти кладуть на
\VUgras{пюпітр} \textsuperscript{(65)}. Один з учнів \VUgras{грає на роялі}
\textsuperscript{(64)}, а \VUgras{учитель музики} \textsuperscript{(66)}
\VUgras{відбиває такт} \textsuperscript{(68)}.

%%K t01-17-1 p
17. --- Ледве видно ріг майстерні \VUgras{ручної праці} \textsuperscript{(69)},
де діти можуть учитися стругати дерево й обпилювати залізо. Не в кожній школі
вона є.

%%K t01-tit-7 sub
\VUsaut{5.0mm}
\VUpk{10.2pt}{40}{Клас природничих наук.}
\VUsaut{3.6mm}

%%K t01-18-1 p
18. --- Учитель математики викладає учням \VUgras{геометрію,} \VUgras{арифметику,
лічбу} і \VUgras{метричну систему}. На таблиці ми бачимо головні геометричні
фігури: \VUgras{коло} \textsuperscript{(a)}, \VUgras{квадрат}
\textsuperscript{(b)}, \VUgras{трикутник} \textsuperscript{(c)},
\VUgras{лінію} \textsuperscript{(d)}.

%%K t01-18-2 p
Ми бачимо також \VUgras{дію}; це не \VUgras{додавання}, не \VUgras{віднімання} і
не \VUgras{множення}: це \VUgras{ділення} \textsuperscript{(e)}.

%%K t01-19-1 p
19. --- На таблиці метричної системи ми бачимо: \VUgras{метр}
\textsuperscript{(a)}, \VUgras{терези} \textsuperscript{(b)} та їхні
\VUgras{важки}, \VUgras{літр} \textsuperscript{(e)}, \VUgras{декалітр}
\textsuperscript{(f)}, \VUgras{децилітр} і \VUgras{кубічний метр}
\textsuperscript{(d)}.

%%K t01-19-2 p
Учні вивчають також \VUgras{природничі науки} \textsuperscript{(11)} ---
зоологію \textsuperscript{(a)}, ботаніку \textsuperscript{(b)}, геологію
\textsuperscript{(c)} --- та \VUgras{історію} \textsuperscript{(12)}.

%%K t01-19-3 p
У шафі стоять прилади для \VUgras{фізики} \textsuperscript{(15)} і для
\VUgras{хімії} \textsuperscript{(16)}.

%%K t01-19-4 p
На шафі стоять: \VUgras{куля} \textsuperscript{(14)}, \VUgras{конус} і
\VUgras{глобус} \textsuperscript{(13)}.

%%K t01-19-5 p
Над таблицями висить \VUgras{мапа}, що зображує п'ять частин світу:
\VUgras{Європу} \textsuperscript{(g)}, \VUgras{Азію} \textsuperscript{(e)},
\VUgras{Африку} \textsuperscript{(f)}, \VUgras{Америку} \textsuperscript{(ab)} і
\VUgras{Океанію} \textsuperscript{(h)}, і два великі океани ---
\VUgras{Тихий} \textsuperscript{(c)} та \VUgras{Атлантичний}
\textsuperscript{(d)}.
\VUsaut{6.0mm}
\VUfilet{22mm}
